\section{Limitations}\label{sec:limitations}

\paragraph{One labeller.} The gold labels are the author's. The second pass measures how
consistent one person is with themselves, not how far two trained annotators would disagree,
so the label ceiling of Section~\ref{sec:results} is optimistic. Gold labels can also be
wrong in ways that favour a model which shares the labeller's reading of the item headings.

\paragraph{The product universe is Open Food Facts.} The gold frame is the Thai-script
products in a crowd-sourced database, not the assortment of Thai shops. It over-represents
packaged, branded products that volunteers chose to enter, and the universe shares $n_k$ in
the simulator are shares of that frame. A statistical office would draw the frame from its
own scanner or web data, where the mix of items, and therefore the error in basis points,
may differ.

\paragraph{Item weights below the published cells are estimated.} TPSO publishes weights for
its aggregates only. The equal split and the NNLS split are two transparent estimates; they
agree on the cells but not on the items inside them. Results that depend on which item
inside a cell carries the weight are reported under both.

\paragraph{The simulator measures classification error only.} It takes official item
relatives as the true price movements of every product and changes nothing but the
assignment of products to items. A real web-price index would add sampling, matching,
quality-adjustment and outlier error on top. Three modelling
choices are assumptions rather than estimates: leaked out-of-scope products move with the
official non-food index, an item that receives no products takes its group's measured
movement, and the product-level dispersion of the noise variants is assumed, not
measured.

\paragraph{Official specifications are not shop listings.} The official labels come from
TPSO's price-collection specifications. They are short and regular;
shop and database texts are longer and noisier. The learning curve therefore
measures the value of this kind of official label, not of labelled shop data.

\paragraph{Closed models change.} The LLMs are reached through a commercial API. Model
revisions, prices and refusal behaviour can change after the run date; every response is
cached so that the numbers here can be reproduced, but not necessarily re-obtained. The
official specifications are public and may have been seen in pre-training. Local compute
is costed at an assumed price per GPU-hour.

\paragraph{Not comparable to Spectrum.} Spectrum classifies euro-area web data into the
ECOICOP 2018 subclasses \citep{bis2026spectrum}. Our taxonomy, language, product universe,
label source and metrics differ, so our numbers are not directly comparable with theirs,
and we do not compare them.
